% document_id: BJP-PROGRAM
% revision: 0.59.0
% as_of: 2026-10-02
% [更新:0.59.0] [確認:0.59.0] — PROGRAM-CORE / REV-590-GUIDE-PROGRAM
% 0.58 review scope: annual raw quarters, climate cache responsibilities and pressure controls; final render and runtime evidence are separate.
% Editable source. tools/build_guides.py --guide PROGRAM_GUIDE
\documentclass[10pt,a4paper,landscape]{article}
\usepackage[margin=12mm,top=15mm,bottom=13mm,headheight=14pt]{geometry}
\usepackage{fontspec,xeCJK}
\setmainfont{Noto Sans CJK JP}\setsansfont{Noto Sans CJK JP}
\setmonofont{Noto Sans Mono CJK JP}[Scale=0.86]
\setCJKmainfont{Noto Sans CJK JP}\setCJKsansfont{Noto Sans CJK JP}\setCJKmonofont{Noto Sans Mono CJK JP}[Scale=0.86]
\usepackage{xcolor}
\definecolor{ink}{HTML}{193247}\definecolor{muted}{HTML}{566A78}
\definecolor{appc}{HTML}{2463A3}\definecolor{weatherc}{HTML}{137D80}\definecolor{dync}{HTML}{B45B25}\definecolor{outc}{HTML}{7451A8}\definecolor{toolc}{HTML}{5C6573}
\usepackage{tikz}\usetikzlibrary{arrows.meta,calc,positioning}
\usepackage{booktabs,tabularx,array,xurl,fancyhdr}
\usepackage[unicode,colorlinks=true,linkcolor=appc,urlcolor=appc]{hyperref}\usepackage{bookmark}
\hypersetup{pdftitle={Balloon-JP プログラム構造ガイド 文書0.59・実装0.58},pdfauthor={Balloon-JP},pdfsubject={実ファイルと関数の包含・依存・呼出・データ受渡し}}
\setlength{\parindent}{0pt}\setlength{\parskip}{3pt}\setlength{\emergencystretch}{2em}
\pagestyle{fancy}\fancyhf{}
\fancyhead[L]{\scriptsize Balloon-JP / プログラム構造ガイド}\fancyhead[R]{\scriptsize 文書0.59 / 実装0.58 / 2026-10-02}
\fancyfoot[L]{\scriptsize BJP-059-PROGRAM / 実構造・操作の受入・科学的採用は別}\fancyfoot[R]{\scriptsize\hyperlink{PROGRAM-MAP}{全体地図へ}\quad\thepage}
\renewcommand{\headrulewidth}{.3pt}
\newcommand{\C}[1]{\texttt{\detokenize{#1}}}\newcommand{\Go}[2]{\hyperlink{#1}{#2}}
\newcommand{\Page}[5]{\clearpage\hypertarget{#2}{}\pdfbookmark[0]{#3}{bm:#2}{\color{#1}\rule{1.4mm}{10mm}}\hspace{3mm}\begin{minipage}[b]{.94\linewidth}{\fontsize{22}{27}\selectfont\bfseries\color{ink}#3}\par{\small\color{muted}#4}\end{minipage}\par\vspace{2mm}{\color{#1}\rule{\linewidth}{.8pt}}\par{\small #5}\par\vspace{1mm}}
\newenvironment{Map}{\begin{tikzpicture}[x=1cm,y=-1cm,>=Latex,every node/.style={font=\small,text=ink},call/.style={->,draw=ink,line width=.75pt},data/.style={->,draw=appc,densely dotted,line width=1.1pt},contain/.style={draw=muted,line width=.7pt},imp/.style={->,draw=muted,dash pattern=on 5pt off 3pt,line width=.85pt}]\path[use as bounding box](0,0) rectangle (27,13.6);}{\end{tikzpicture}}
\newcommand{\Card}[7]{\node[anchor=north west,draw=#2,fill=#2!6,rounded corners=1.6mm,line width=.65pt,text width=#5cm,inner sep=2.7mm,align=left] (#1) at (#3,#4) {\textbf{\color{#2}#6}\par\vspace{1mm}#7};}
\newcommand{\File}[7]{\node[anchor=north west,draw=#2,fill=#2!5,line width=.85pt,text width=#5cm,inner sep=2.7mm,align=left] (#1) at (#3,#4) {\textbf{\color{#2}#6}\par\vspace{1mm}#7};}
\newcommand{\Tip}[7]{\node[anchor=north west,draw=#2!70,fill=white,rounded corners=2mm,dashed,text width=#5cm,inner sep=2.7mm,align=left] (#1) at (#3,#4) {\textbf{\color{#2}#6}\par\vspace{1mm}#7};}
\newcommand{\Small}[4]{\node[anchor=north west,text width=#3cm,align=left,font=\footnotesize,text=muted] at (#1,#2) {#4};}
\catcode`\_=12

\begin{document}
% LOGIC-BEGIN:PROGRAM-CORE
\Page{appc}{PROGRAM-READ}{変更するために、構造を読む}{利用者の比較 → 画面・サービス・核 → 担当ファイル → 主要な機能}{この資料は実装の地図である。式の導出や細かな入力仕様は専用資料へ送り、どこを読むか・変更がどこへ届くかを短い経路で示す。}
\begin{Map}
\Card{whole}{ink}{.2}{.4}{7.3}{全体の問い}{\Go{PROGRAM-MAP}{何がどこにあるか}\par \Go{PROGRAM-DEPENDENCY}{誰が誰を使うか}\par 草案、固定結果、表示、計算、保存を見分ける。}
\Card{branch}{weatherc}{9.4}{.4}{7.2}{一つの枝へ降りる}{\Go{WEB-FRONTEND}{画面} ／ \Go{WEB-SERVICE}{サービス}\par \Go{WEATHER-FILE}{場} ／ \Go{WEB-CLIMATE-BACKEND}{風統計} ／ \Go{DYNAMICS-FILE}{一飛行}\par \Go{ENSEMBLE-CORE}{標本集合} ／ \Go{CLI-EXPORT}{出力}\par 枝の責務、入出力、一段下の実体を見る。}
\Card{change}{dync}{18.6}{.4}{7.2}{変更と検証へ戻る}{\Go{PROGRAM-CHANGE}{変更の逆引き}\par 親の目的と横の契約を戻り先にする。全関数を同時に覚えなくてよい。}
\draw[call](whole.east)--(branch.west);\draw[call](branch.east)--(change.west);
\Card{colors}{ink}{.2}{4}{12}{役割を色で追う}{\textcolor{appc}{青：操作の入口}　\textcolor{weatherc}{緑：環境場}\par \textcolor{dync}{橙：飛行・物理・積分}　\textcolor{outc}{紫：結果と保存}\par \textcolor{toolc}{灰：試験・資料・管理}\par 色だけに意味を持たせず、実パスと責務を併記する。}
\Card{lines}{ink}{14.1}{4}{11.7}{線と箱を読み分ける}{細線：実体の包含。実線矢印：主な呼出。\par 灰破線：import依存。青点線：値の受渡し。\par 四角：実ファイル/クラス。丸角：機能。\par 図は静的な全呼出網ではなく、その枝を理解するための主要な関係を選ぶ。}
\Tip{basis}{appc}{.2}{8.1}{12}{実装と互換入口を分ける}{0.24で責務を実パッケージへ分けた。旧\C{weather.py}/\C{dynamics.py}は公開名を保つ再export入口で、別実装を持たない。未実装の慣性・熱・気象誤差を、現役の仮感度集合へ混ぜない。}
\Tip{books}{outc}{14.1}{8.1}{11.7}{地図の次に読む資料}{\path{IMPLEMENTATION_NOTES.md}：処理/入出力の詳細\par \path{THEORY_GUIDE.pdf}：式・仮定・数値法\par \path{COMMANDS.md}：実行と再現\par \path{IMPLEMENTATION_INDEX.json}：実定義・公開名・静的依存}
\Small{.2}{12.1}{26}{文書改訂0.59・実装0.58。固定PR31 C2 95f09d2fc0aeを基準に、気象画面の担当と参照先を局所訂正した。コードと科学式は変更しない。0.58の年間48区分、年間値の再利用、気圧面の操作枠は継承する。全枝の科学的再検証ではなく、既存の受入・未確認はS36へ戻る。式はTHEORY、mainへの統合は別判断とする。}
\end{Map}

\Page{ink}{PROGRAM-MAP}{比較する操作を、実際の枝へ結ぶ}{階層0 / repo → frontend・backend・balloon_sim}{草案を編集しても固定結果を変えない。そのため、画面の選択、計算の受入、物理の意味を別の実体に置く。}
\begin{Map}
\File{root}{ink}{.2}{.2}{25.6}{private repo：同じcommitから起動する正本}{\C{frontend/} と \C{backend/} を通常のGit差分で育て、既存 \C{balloon_sim/} と保存fixtureを同じ版へ結ぶ。}
\File{fe}{appc}{.2}{3}{7.2}{\Go{WEB-FRONTEND}{\C{frontend/}}}{React/TypeScriptと継承した画面素材。草案・固定結果の選択、地図、履歴、気象画面。Viteが開発/ビルドを担う。}
\File{be}{outc}{9.4}{3}{7.2}{\Go{WEB-SERVICE}{\C{backend/}}}{FastAPIとapplication service。\par 固定入力・一飛行・結果・草案。\par \Go{WEB-WEATHER}{weather/}：取得。\Go{WEB-ENSEMBLE-BACKEND}{ensemble/}：固定集合。\Go{WEB-CLIMATE-BACKEND}{climate/}：平均・原標本の風統計。}
\File{kernel}{dync}{18.6}{3}{7.2}{\C{balloon_sim/}}{\Go{WEATHER-FILE}{environment}：保存場\par \Go{DYNAMICS-FILE}{flight}：一飛行\par \Go{ENSEMBLE-CORE}{ensemble}：抽選・共通集計\par \Go{CLI-EXPORT}{results}：一飛行の出力。}
\draw[contain](root.south)--++(0,.3)-|(fe.north);\draw[contain](root.south)--(be.north);\draw[contain](root.south)--++(0,.3)-|(kernel.north);
\File{support}{toolc}{.2}{7}{7.2}{\C{tests / docs / tools}}{守る契約、構造/理論/操作の説明、生成/点検。計算の実行依存へ混ぜない。}
\File{inputs}{weatherc}{9.4}{7}{7.2}{\C{references / examples}}{保存GFS・JRA原本と機体入力の例。現在予報や実機係数の採用値として扱わない。}
\File{state}{outc}{18.6}{7}{7.2}{repo外の利用状態}{草案/仕事・完成場・固定結果は\C{BALLOON_DATA_DIR}へ。外生DB/原標本manifestの読取元は起動設定で別に指定。Gitソースや\C{dist}と別の寿命。}
\Tip{why}{ink}{.2}{10.5}{25.6}{この分担が守ること}{画面を編集しても保存結果の入力が変わらず、HTTPやDBを替えても物理核へ依存を逆流させない。表示・入力・場・数値の意味は別の枝へ戻れる。フォルダーを増やすことが目的ではない。}
\end{Map}

\Page{appc}{PROGRAM-DEPENDENCY}{取得・飛行・風統計を、別の枝でつなぐ}{階層1 / frontend → application → weather・ensemble・climate・一飛行worker}{集合は既存の一飛行を同じqueueへ投入する。風統計は独立した集約/保存serviceへ渡し、飛行workerへは送らない。}
\begin{Map}
\File{fe}{appc}{.2}{.3}{7.2}{\C{frontend/src/}}{草案・固定結果を別に保持。\par \C{api.ts}を共用し、\C{ensembleApi.ts / climateApi.ts}から各APIへ。}
\File{service}{outc}{9.4}{.3}{7.2}{\C{app.py → application.py}}{HTTP → ApplicationService。\par 保存GFS/JRAと完成取得場を一覧化し、入力を固定。\par 既存queueから一workerへ送る。}
\File{worker}{dync}{18.6}{.3}{7.2}{\C{backend/worker.py}}{\C{execute_run}：spawnした一processで保存場 → simulate → export。集合内の各試行も同じ一飛行。}
\draw[data](fe.east)--(service.west);\draw[call](service.east)--(worker.west);
\File{acquire}{weatherc}{.2}{4.4}{7.2}{\Go{WEB-WEATHER}{\C{weather/service.py}}}{取得job専用thread。原本から完成場を作る。候補への適用は画面で明示する。HTTPは共有gateway。}
\File{ensemble}{outc}{9.4}{4.4}{7.2}{\Go{WEB-ENSEMBLE-BACKEND}{\C{ensemble/service.py}}}{固定計画と全試行台帳。\par \C{tick}が一飛行を既存queueへ。\par 完了結果を固定snapshotと分析へ。}
\File{core}{dync}{18.6}{4.4}{7.2}{\C{balloon_sim/}}{environment：保存場の再構成。\par flight：一飛行。\par \Go{ENSEMBLE-CORE}{ensemble}：抽選と集計。\par HTTP・DBを持たない。}
\draw[call](service.south)--++(0,.45)-|(acquire.north);\draw[call](service.south)--(ensemble.north);\draw[call](worker.south)--(core.north);
\draw[imp](ensemble.east)--node[above,font=\scriptsize]{抽選・集計}(core.west);
\File{climate}{weatherc}{.2}{9.1}{12}{\Go{WEB-CLIMATE-BACKEND}{\C{climate/service.py}}}{ApplicationServiceが起動し、HTTPから直接呼ぶ。\par 独立lock/SQLiteで集約結果を固定。DBは\C{source.py}、原標本は\C{archive.py}へ。飛行共通lockを占有しない。}\draw[call](service.south west)--++(-.55,.35)--(8.85,8.65)-|(climate.north);
\Tip{boundary}{outc}{14.1}{9.1}{11.7}{固定物の種類を分ける}{完成場 → 一飛行の入力。\par run結果 → 集合snapshot → 選択群分析。\par 風統計artifact → 年間/地域/季節図。\par 集約統計を瞬時の飛行場へ変換しない。}
\end{Map}

\Page{appc}{WEB-FRONTEND}{frontend：共通の計画と、用途ごとの画面}{階層2 / frontend → src → App・状態・screens}{育てた予報・詳細・気象・季節の画面を、同じ計画と固定結果へ結ぶ。Reactの共通状態と、各画面の描画状態は異なる枝にある。}
\begin{Map}
\File{root}{appc}{.2}{.2}{25.6}{\C{frontend/src/}：\C{main.tsx}から起動}{\C{main.tsx}がReactとLeaflet CSSを読み込む。\C{App.tsx}が共通ヘッダー・用途切替・保存入口を構成し、\C{platform.css}がその配置を担当する。}
\File{app}{appc}{.2}{2.35}{7.2}{\C{App.tsx}}{計画/気象/季節の行先を選ぶ。\par 候補の追加・延期・独立化を共通計画へ反映。画面の変更通知を受け、用途別の表示状態を保存側へ渡す。}
\File{state}{outc}{9.4}{2.35}{7.2}{\C{useWorkspace.ts}\\\C{workspaceState.ts}}{草案・固定run・選択・仕事・保存。\par \Go{WEB-RECOVERY}{固定要求・復旧の枝へ}。\par \C{jobObservation.ts}へ単便のGET監視を委譲。取得UIも共有する。\par 集合は\Go{WEB-ENSEMBLE-FRONTEND}{useEnsembleWorkspace}へ。表示選択tokenを共有。}
\File{api}{outc}{18.6}{2.35}{7.2}{\C{api.ts}}{相対\C{/api/v1}のJSON送受信。\par \C{readRequest}：期限付きGET。\par \C{request}：一般送受信。集合/風統計の専用APIも利用する。\Go{WEB-SERVICE}{backend}へ。}
\draw[contain](root.south)--++(0,.18)-|(app.north);
\draw[call](app.east)--(state.west);\draw[call](state.east)--(api.west);
\File{forecast}{appc}{.2}{6.15}{12}{\C{screens/ForecastScreen.tsx}}{予報・詳細・季節計画を同じ画面素材で開く。\par 実結果と現在sourceをprojectionへ渡しruntimeへ同期。未読時も保存した群/時刻を保持。専用portalへ\C{ConfigEditor}・\C{SamplingEditor}・固定結果選択・\Go{WEB-PREPARATION}{\C{WeatherPreparation}}を載せる。\par 取得準備と入力編集は専用の枝へ降りる。}
\File{weather}{weatherc}{14.1}{6.15}{11.7}{\C{screens/WeatherScreen.tsx}}{人工の四画面と、実風統計の年間/地域/季節図を選ぶ。\par 実側の\C{useClimateWorkspace}が草案query/固定artifact、runtimeが表示を担う。\par 人工の提案を保持し、実側は\Go{WEB-HISTORICAL}{原日時の関心}をAppへ。\Go{WEB-CLIMATE-FRONTEND}{統計と共通図の枝へ}。}
\draw[data,<->](app.south)--node[left,font=\scriptsize]{条件・結果 / 操作}(forecast.north);
\draw[data,<->](app.south)--++(0,.55)-|(weather.north);
\node[anchor=south,font=\scriptsize,fill=white,inner sep=1pt] at (16.3,5.83) {表示状態 / 季節への引渡し};
\Tip{types}{dync}{.2}{10.05}{12}{横で共有する型・入力}{\C{domain.ts}：草案・固定結果、気象の内容照合、JST/UTCの秒を保つ変換。\par \C{fixedResultRead.ts}：読取状態・表示名・詳細待機の判定。各hookと画面が共有し、HTTPを持たない。\par \C{ConfigEditor.tsx}：草案入力。固定結果は不変。}
\Tip{ownership}{appc}{14.1}{10.05}{11.7}{もう一段下へ降りる}{\Go{WEB-SCREENS}{runtimeとcontrollerの枝へ}。\par 画面の共通時刻・群選択・地図視点は各controller側。Appが全描画状態を直接操作する構成ではない。境界の詳細は実装ノートへ。}
\end{Map}

\Page{appc}{WEB-SCREENS}{screens：画面素材を、限定した描画領域へ載せる}{階層3 / screens → forecast・weather・shared・lifecycle}{各React画面がhostを持ち、runtimeが既存の画面と描画処理を載せる。画面の内容を一から描き直した構成ではなく、具体素材を取り込んだ移行途中の構成である。}
\begin{Map}
\File{runtime}{appc}{.2}{.3}{7.2}{\C{forecast/runtime.ts}\\\C{weather/runtime.ts}}{markup/CSSを読み、Leaflet・Plotly等と補助モジュールを渡してcontrollerを起動。終了時は描画の破棄処理へつなぐ。}
\File{controller}{appc}{9.4}{.3}{7.2}{\C{forecast/controller.js}\\\C{weather/controller.js}}{前者は候補/群/詳細・履歴・3D連絡と地図の初回fit/視点復元。後者は人工画面を保持し、実統計はclimateControllerへ委譲。画面内DOMと操作を主に所有する。}
\File{scope}{toolc}{18.6}{.3}{7.2}{\C{lifecycle.ts}}{runtimeが\C{screenScope}を生成。host中心のdocument/window、イベント・timer・frame・observerの寿命を束ねる。\par controllerから終了時にdispose。}
\draw[call](runtime.east)--(controller.west);\draw[call](controller.east)--node[above,font=\scriptsize]{終了処理}(scope.west);
\File{source}{appc}{.2}{4.7}{7.2}{各枝の\C{markup.html}\\\C{screen.css}}{既存画面のDOM配置と画面単位のCSS。\par Reactはhostを所有し、controllerがその内部を描画。実入力の専用portalだけReactが担当する。}
\File{projection}{outc}{9.4}{4.7}{7.2}{\C{forecast/projection.ts}}{n=1をタブ・標本・履歴へ変換。\par 集合は\Go{WEB-ENSEMBLE-FRONTEND}{ensembleProjection.ts}へ分岐。草案差分と気象の内容照合を分け、計算時の支持境界を保持。全台帳・地図座標・遅延読込履歴を別にする。}
\File{helpers}{weatherc}{18.6}{4.7}{7.2}{\C{forecast/}・\C{weather/}\\\C{shared/}の補助実体}{forecast：図形/外観・名称・保存・scene文書。weather：資料別adapterと共通図・保存、人工fixture・風向・影響地図。shared：色/帯/名称の表現。旧補助はruntimeが注入。新しい統計/共通図は直接import。}
\draw[imp](runtime.south)--(source.north);
\draw[data](projection.north)--node[left,font=\scriptsize]{画面経由で同期}(controller.south);
\draw[data](helpers.north)--++(0,-.45)-|(controller.east);
\node[font=\scriptsize,fill=white,inner sep=1pt] at (21.7,3.95) {直接import / 旧補助をruntime注入};
\Tip{scene}{outc}{.2}{9.05}{12}{\Go{WEB-SCENE3D}{3Dの表示と読取りの枝へ}}{controllerが\C{scene-document.js}で対象の固定snapshotを作り、\C{postMessage}で\C{/scene3d/index.html}へ渡す。表示設定は親へ戻し、KMLと保存に共有。配信元は\C{frontend/scene3d/}。\C{src/main.js}が地表overlay・空中route・hoverの寿命を束ねる。ViteがnpmのCesiumと固定高度格子をdistへ置く。}
\Tip{debt}{toolc}{14.1}{9.05}{11.7}{現在の構造で残る制約}{大きなcontrollerに複数の操作・描画責務が残り、型のないbridgeやscopeを介す。移行だけで分割完了とはしない。旧\C{MapPanel.tsx}/\C{HistoryPanel.tsx}は残るが、現在Appの表示経路からは呼ばれていない。}
\Small{.2}{12.8}{26}{図は読取時点の実ソースの関係。人工表示、実n=1/仮感度集合、実DBの期間統計、保存JRAの元日時場を区別する。任意過去窓の取得UI・代表性を定めた多年季節MCは未実装。登録済み原日時の有限比較は\Go{WEB-HISTORICAL}{原日時の枝}へ。controllerの寿命・操作全体・写真3D接続の受入を、この配置図だけで認定しない。}
\end{Map}


\Page{outc}{WEB-SCENE3D}{scene3d：同じ経路を、地形との関係へ開く}{階層4 / frontend → scene3d → src と data}{写真地図の接続、固定結果の図形、地物の読取りは別の寿命を持つ。空中線を追加しても地表ホバーを壊さず、表示変更だけでGoogleへ接続し直さない。}
\begin{Map}
\File{host}{appc}{.2}{.1}{25.6}{\C{index.html → src/main.js}：表示のhost}{親のforecast/controllerから固定snapshotと表示設定を受ける。provider.jsへ明示接続を渡し、候補catalog、地表overlay、空中route、視点・読取りを束ねる。表示の変更通知は親の保存UI/KMLへ戻す。}
\File{contract}{outc}{.2}{2.65}{7.2}{\C{src/protocol.js}}{候補/結果/選択群とgeometryを検査。\par 地表線・空中線・幕の経緯度/相を対応づけ、ASL/人工AGLを区別する。\par 生成元は親のscene-document。\par allObjectsはKML用、3Dは可視対象。}
\File{ground}{outc}{9.4}{2.65}{7.2}{地表の表示}{\C{src/overlays.js}\par 点・帯・領域・地表線のCesium定義。\par \C{src/overlay-store.js}\par 新定義を準備して差替え、失敗時は旧sourceへ戻す。}
\File{route}{outc}{18.6}{2.65}{7.2}{\C{src/route-layer.js}}{空中表示と高度準備の寿命。\par 実ASLは地物照会なしで全経路。人工AGLだけ有限照会・取消を持つ。\par 接続/視点は生成しない。\par checkpointが旧高度と表示を保持。}
\draw[contain](host.south)--++(0,.3)-|(contract.north);\draw[contain](host.south)--(ground.north);\draw[contain](host.south)--++(0,.3)-|(route.north);
\File{read}{appc}{.2}{7.05}{7.2}{地物を読む補助群}{\C{surface-session / surface-pick}\par 読取り順と通常描画待ち。\par \C{hover-controller / hover-labels}\par 静止と名称、帯はband-readout。\par 表示中/退役待ち図形を地物から除外。}
\File{geometry}{outc}{9.4}{7.05}{7.2}{\C{src/route-geometry.js}}{原経路の間に表示節点を追加。\par 同じ節点と曲率から空中線・連続幕を作る。\par 幕の地下余長は写真meshに隠れる。\par 空中要素は透過passとし、背後の地物深度を読めるようにする。}
\File{datum}{weatherc}{18.6}{7.05}{7.2}{\C{src/vertical-datum.js}}{表示専用のEGM96格子補間。\par \C{data/WW15MGH.DAC}\par 固定15分格子。data/READMEに出典。\par route-layerがローカルfetch。\par 科学ASL・KML値を変更しない。}
\draw[imp](route.south)--++(0,.4)-|(geometry.north);\draw[imp](route.south)--(datum.north);
\Tip{boundary}{ink}{.2}{11.15}{25.6}{どこへ戻って判断するか}{高度の意味と表示限界はTHEORY、切替・再接続・保留の操作はCOMMANDSへ。幕は上下の対応を示す。地形高の取得可否で実経路を分断せず、描画の重なりを衝突解析へ読み替えない。新しい高さ表示のために飛行核や保存結果のschemaを変更しない。視点を動かす操作はmainとcamera-fit/view-memoryへ戻る。}
\end{Map}

\Page{weatherc}{WEB-CLIMATE-FRONTEND}{weather：資料を分け、図と保存を共有する}{階層3 / WeatherScreen → 状態・資料adapter → 図別controller・plots・export}{人工標本と実集約値では読める量が違う。実側は固定結果内の区切りを選ぶ。原標本の48区分は年間専用。月中央値/幅、補助の平均、代表点風配はそれぞれの意味を保つ。}
\begin{Map}
\File{host}{appc}{.2}{.2}{25.6}{\C{screens/WeatherScreen.tsx}：\C{FixtureWeatherScreen / ClimateWeatherScreen}}{人工側は既存runtime/状態を接続。実側はhookのartifactを粒度別VMへ変換し、Reactのquery入力と\C{runtime.sync}をつなぐ。用途切替と保存計画はAppと共用。関心の受渡しは\Go{WEB-HISTORICAL}{\C{ClimateFlightHandoff.tsx}}へ。}
\File{state}{outc}{.2}{3}{7.2}{\C{useClimateWorkspace.ts}}{草案queryと固定artifactを分ける。気圧面/風配点の探索を同じ最新要求列へ。\par 再送ID・応答寿命・GET復帰を保持。入力/支持は\C{climateDomain.ts}。参照はweather_realへ保存。}
\File{api}{outc}{9.4}{3}{7.2}{\C{climateApi.ts → api.ts}}{sources / create / get。\par \Go{WEB-CLIMATE-BACKEND}{固定統計のAPI}へJSON送受信。\par 集計条件の適用はPOST、保存図のGET再読は再集計しない。}
\File{display}{appc}{18.6}{3}{7.2}{\C{weather/climateController.js}}{図別区切りと元格子点を選ぶ。年間だけ48区分を追加し、保存図の切替はPOSTなし。\par 面操作はslider・選択/表示値・状況の枠を分ける。人工側はcontroller.js。}
\draw[contain](host.south)--++(0,.35)-|(state.north);\draw[call](state.east)--(api.west);\draw[contain](host.south)--++(0,.35)-|(display.north);
\File{adapters}{weatherc}{.2}{7}{12}{\C{weather/sources/}}{\C{fixtureSummary.ts}：人工原標本のaggregateFixture → VM。\par \C{climateSummary.ts}：固定summary → 粒度別VM。年間48区分の暦・件数を独立検査。\par 原標本分位の検査は\C{../monthlyProfiles.ts}を参照。旧artifactの図も保つ。}
\File{vm}{weatherc}{14.1}{7}{11.7}{\C{weather/viewModel.ts}}{\C{WindViewModel}：時期・面・格子、平均量、母数、定義と固定来歴。分位と風配は別の値/count/階級を保持。\par \C{levelEdges / niceSpeedMax}：表示区画/尺度。}
\draw[data](state.south)--++(0,.5)-|node[pos=.65,above,fill=white,inner sep=1pt,font=\scriptsize]{固定summary}(adapters.north);\draw[imp](adapters.east)--(vm.west);
\File{plots}{weatherc}{.2}{10.4}{12}{\C{weather/}：図の意味ごとの担当}{\C{plots.ts}：年間/地域。\C{climateDistributions.ts}：平均/風配・静穏と元格子の選択図。\C{monthlyProfiles.ts}：固定12か月の中央値/幅・共通軸。平均から分位を作らない。}
\File{export}{outc}{14.1}{10.4}{11.7}{\C{weather/export.ts}}{保存開始時の図/条件/来歴を固定し背景を埋める。\par 年間/地域・平均/風配・月分位の図を保存。}
\draw[data](adapters.south)--(plots.north);\draw[imp](export.west)--(plots.east);
\end{Map}

\Page{appc}{WEB-PREPARATION}{気象の準備：保存場の選択と、GFSの取得}{階層3 / ForecastScreen → 保存場・取得UI・入力の操作}{同じ画面の共通準備と条件棚を使う。取得完了で増えるのは利用可能な場であり、草案や表示中の固定結果を自動で置換しない。}
\begin{Map}
\File{host}{appc}{.2}{.2}{25.6}{\C{screens/ForecastScreen.tsx}：既存画面とReactの接点}{\C{#realWeather} portalへ取得準備を載せる。\C{bridge.acquisitionPreview()}でcontrollerへ丸め後の橙破線境界・run・必要窓・古化を渡す。固定結果の灰色支持境界は別に保持する。}
\File{ui}{appc}{.2}{3}{7.2}{\C{WeatherPreparation.tsx}}{保存例・部分一覧・計画・明示適用。\par \C{jobObservation.ts}へ取得監視を委譲。単便と共有する。\par 完了時はsourcesだけ再取得。\par plan/jobの正本はbackend。}
\File{rules}{weatherc}{9.4}{3}{7.2}{\C{weatherAcquisition.ts}}{API型、親子の必要時刻窓、取得要求、run/leadのvalid表示。\par 適用前の時間・地点を確認。\par 全飛行の支持を保証しない。}
\File{api}{outc}{18.6}{3}{7.2}{\C{api.ts / useWorkspace.ts}}{api：inventory/plan/job/sourceのHTTP。\par useWorkspace：既存草案・run・保存、\C{refreshSources}。取得結果で草案を作り直さない。}
\draw[contain](host.south)--++(0,.35)-|(ui.north);\draw[call](ui.east)--(rules.west);\draw[call](ui.north east)--++(0,-.45)-|node[pos=.75,above,font=\scriptsize]{HTTP操作}(api.north);
\File{edit}{dync}{.2}{7}{7.2}{\C{ConfigEditor.tsx}\\\C{modelDrafts.ts}}{JST秒入力をUTCへ変換。\par \C{switchModel}が方式別値を退避/復元。非選択値は送信configへ混入させない。地表の明示適用を草案へ返す。}
\File{ground}{weatherc}{9.4}{7}{7.2}{\C{GroundPreparation.tsx}}{場・地点・時刻を固定して\C{api.ground}へ照会。\par 後着応答の識別、地表との差、指定余裕の明示適用。別の地表式は持たない。}
\File{apply}{outc}{18.6}{7}{7.2}{\C{App.tsx → modelDrafts.ts}}{\C{applyCandidateWeather}が親子のsourceを更新し、物理入力を保持。新候補の作成だけ入力例を複写。共通projectへ反映。}
\draw[call](edit.east)--node[above,font=\scriptsize]{照会UI}(ground.west);
\draw[call](ground.north)--++(0,-.35)-|(api.south);\node[font=\scriptsize,anchor=south] at (18,6.6) {GET};
\draw[data](ui.west)--++(-.35,0)--(-.15,10.65)-|(apply.south);\node[font=\scriptsize,anchor=south] at (15,10.6) {気象の明示apply};
\Tip{state}{outc}{.2}{11.1}{25.6}{保存先と、別の寿命を持つもの}{\C{ui_state.forecast_real.weatherPreparation}：フォーム、手順位置、plan/job/request参照、preview。\C{editor_models}：非選択方式の編集値。固定runは別参照。実風統計は別のweather_real状態と固定artifactへ。原日時試算は\Go{WEB-HISTORICAL}{season_realの別状態}へ。取得GFSや半月統計を原日時場へ読み替えない。}
\end{Map}

\Page{appc}{WEB-ENSEMBLE-FRONTEND}{集合の画面：固定結果を、選んで掘り下げる}{階層3 / ForecastScreen → 入力・非同期状態・描画adapter}{既存の候補・延期列・地図・詳細画面へ集合を接続する。仕事の進捗と表示中のsnapshotを別に保持し、完了を自動表示しない。}
\begin{Map}
\File{editor}{appc}{.2}{.3}{7.2}{\C{SamplingEditor.tsx}}{仮幅/理由/N/seed、対応候補。\par 固定計画の確認 → 明示実行。\par \C{ensembleDomain.ts}が型・元量と方式・親子計画・固定参照を束ねる。}
\File{state}{outc}{9.4}{.3}{7.2}{\C{useEnsembleWorkspace.ts}}{Appが呼び、共通画面の寿命で保持。\par 計画・job・case・分析・履歴のcacheと寿命。\par case読取状態とGET再読。内部の\C{createEnsembleObserver}がGET観測/失敗後の待機を管理。\par \C{ensembleApi.ts}でGET期限/abortを接続。}
\File{host}{appc}{18.6}{.3}{7.2}{\C{ForecastScreen.tsx}}{既存portalへ入力を載せる。\par 固定caseを\C{collectionResult}で変換し、現在sourceとの照合は別に加えてタブへ同期。\par bridgeで群分析/個別履歴を要求。}
\draw[call](editor.east)--(state.west);\draw[data](state.east)--(host.west);
\File{projection}{outc}{.2}{5.05}{12}{\C{forecast/ensembleProjection.ts}}{\C{collectionResult}：全台帳と地図点を分離。\par 気象の由来照合は\C{domain.ts}と共用。\par \C{empiricalContours / analysisRoutes / aggregateTraces}：共通分析の幾何・相別線/帯を描画値へ。\par \C{pairedRows}：draw_idで対応を結ぶ。}
\File{controller}{appc}{14.1}{5.05}{11.7}{\C{forecast/controller.js → math.js}}{controller：標本選択・図・3D連絡。mathは領域所属判定。\par \C{analysisLoadState.ts}：全ID一致の既存分析を再利用し、部分群の待機/失敗を保持。未送信だけ担当取得後に再開。\par \C{historyLoadQueue.ts}：集計から独立した原履歴GETを画面内4並列で読み、未開始の旧選択を破棄する。}
\draw[call](host.south)--++(0,1.1)-|(projection.north);\draw[data](projection.east)--(controller.west);
\Tip{saved}{outc}{.2}{10}{12}{保存する参照、別に残る状態}{\C{Project.compare_results}：固定run/case。\par \C{ui_state.forecast_real.selectionContext}：群の分析ID/key。再読はgetAnalysisへ。\par 計画保存は\C{useWorkspace}が担当する。}
\Tip{limit}{ink}{14.1}{10}{11.7}{実装上の境界}{地図点がない停止を全台帳から消さない。\par 禁止域の所属はfrontend判定であり、backendによる幾何の独立再検証ではない。\par rendererは既存2D/3Dを使う。}
\end{Map}

\Page{outc}{WEB-SERVICE}{backend：仕事と固定結果の寿命を管理する}{階層2 / backend → HTTP・application・worker・storage・weather・ensemble・climate}{ApplicationServiceは単一利用者・一保存先を所有する。計算結果と完成した気象assetを、途中のdirectoryから区別して受け渡す。}
\begin{Map}
\File{http}{appc}{.2}{.3}{7.2}{\C{app.py / contracts.py}}{create_app、JSON外枠、API/配信。\par \C{errors.ServiceError}をHTTPへ。\par 物理検証はvalidate_config。取得/集合/風統計の要求は各枝の契約へ。}
\File{application}{outc}{9.4}{.3}{7.2}{\C{application.py}}{\C{ApplicationService}\par 保存GFS/JRAを登録し、草案・run・queueを管理。\C{weather_ground}は既存場の地表を読む。\par weather/ensemble/climateの親。集合はDB/workerを共有、風統計は\Go{WEB-CLIMATE-BACKEND}{別lock/SQLite}。}
\File{worker}{dync}{18.6}{.3}{7.2}{\C{worker.py}}{固定入力を照合し、\C{load_weather → simulate → export_result}へ。構築した場の実効metadataを出力。source_snapshotはclimateも含むが、風統計の集約は行わない。}
\draw[call](http.east)--(application.west);\draw[call](application.east)--(worker.west);
\File{db}{outc}{.2}{4.8}{7.2}{SQLite / \C{registry.sqlite3}}{草案revision、固定run仕様、可変job状態、固定結果への参照。同じ要求IDの再送は同じrunへ戻す。}
\File{store}{outc}{9.4}{4.8}{7.2}{\C{storage.py}}{\C{specs / staging / results}\par 核の7出力を照合し、run仕様とCOMMITTEDを加えrename。読込時もhashを検査。}
\File{source}{weatherc}{18.6}{4.8}{7.2}{\Go{WEB-WEATHER}{\C{weather/}：取得と完成場}}{完成取得assetをApplicationServiceへ。\par 保存GFS/JRA例との一覧化とrunへの固定は親の責務。\par 取得計画・job・原本の保存は、この枝へ降りて読む。}
\draw[call](application.south)--(store.north);\draw[data](store.west)--node[above,font=\scriptsize]{結果参照}(db.east);\draw[data](source.north)--++(0,-.8)-|node[pos=.65,right,font=\scriptsize,fill=white,inner sep=1pt]{ID/場来歴}(application.east);
\Tip{limit}{ink}{.2}{9.7}{25.6}{仕事の種類と、読取りの経路を分ける}{飛行jobは待機中だけ取消、実行中は完了を待つ。取得は協調取消と明示retry。集合は\Go{WEB-ENSEMBLE-BACKEND}{ensemble/}へ。地表照会は\C{app → application → load_weather → ground_altitude}の読取経路で、飛行queueへ送らない。起動source変更と旧結果の閲覧は別。風統計はHTTPから独立serviceへ渡す。}
\end{Map}

\Page{outc}{WEB-ENSEMBLE-BACKEND}{ensemble：全試行台帳から、不変の集合へ}{階層3 / backend → ensemble → 契約・計画・service・保存・表示用台帳}{ApplicationServiceのDBとlockを共有するapplicationの枝。並列計算基盤を追加せず、既存の一飛行を同じqueueへ投入する。}
\begin{Map}
\File{contract}{appc}{.2}{.3}{7.2}{\C{contracts.py}}{Sampling / Case / PlanRequest、\par HistoricalPlanRequest、\par Submit / Cancel / Retry / AnalysisRequest。\par 有限入力と参照を受け入れ、計算/保存をここで行わない。}
\File{plan}{dync}{9.4}{.3}{7.2}{\C{planning.py / build_plan}}{同じ元実現値を候補/延期へ展開。\par source・場・全時間窓・運用予算を固定し、不成立入力も残す。\par 抽選は\Go{ENSEMBLE-CORE}{核sampling}へ。\par 原日時は\Go{WEB-HISTORICAL}{historical_planning.py}へ。}
\File{views}{outc}{18.6}{.3}{7.2}{\C{views.py / counts・compact}}{全予定数と状態別数を返す。\par trial/draw/attemptの同一性と終点の短い情報を作る。\par raw履歴とは別の台帳表示。}
\File{service}{outc}{.2}{5}{12}{\C{service.py / EnsembleService}}{plan → submit → tick → 完了の回収 → snapshot。\par tickは\C{ApplicationService._register_fixed_run}へ。\par cancelは未開始を止める。retryだけ同じ元標本の新attemptへ。\par analyzeは固定case/選択IDを核statisticsへ渡す。}
\File{store}{outc}{14.1}{5}{11.7}{\C{storage.py / publish・read}}{SQLite：計画・集合・trial・attempt・固定参照。\par \C{ensemble-artifacts/}：計画/snapshot/分析。\par documentとmarkerをpartialへ書いてrename。\par read時もhashを検査する。}
\draw[imp](service.north)--(contract.south);\draw[call](service.north east)--++(0,-.5)-|(plan.south);\draw[call](service.north east)--++(0,-.5)-|(views.south);\draw[call](service.east)--(store.west);
\Tip{core}{dync}{.2}{10}{12}{外へ委譲する二つの計算}{抽選/集計：\C{balloon_sim/ensemble}。\par 一飛行：既存worker → flight/results。\par 履歴byte予算を超えた分析は理由を返し、着地点と全台帳を保持。再開で旧snapshotを変えない。}
\Tip{scope}{ink}{14.1}{10}{11.7}{現在の境界と戻り先}{群は選択ID・領域・分類器版と共に保存。禁止域の所属はfrontend側で、サーバは独立再判定しない。数式はTHEORY、API/取消/復旧の詳細は実装ノートへ。}
\end{Map}


\Page{outc}{WEB-RECOVERY}{受付と保存先：不明を別の実行へ変えない}{階層3 / frontendの要求台帳 → backendのinstance・固定spec・完成結果}{入力草案、送信した要求、受理済みの仕事、表示結果の寿命を分ける。同じURLであることだけでは、同じ保存台帳とは分からない。}
\begin{Map}
\File{ledger}{outc}{.2}{.3}{7.2}{\C{submissionLedger.ts}}{固定要求のlocalStorage台帳。\par 完全本文・ID・接続先を送信前に保持。不正な記録は再送を止めて保つ。\par put/removeで送信権を確認。結果本体やAPIキーは置かない。}
\File{single}{appc}{9.4}{.3}{7.2}{\C{singleSubmission.ts}\\\C{useWorkspace.ts}}{親/延期列を一件ずつ投入。\par 不明なら読取照合と明示続行。保存PUTは送信草案と後続編集を分離。\par 単便は\C{api.ts}へ。\C{jobObservation}はGETだけ。}
\File{mutation}{appc}{18.6}{.3}{7.2}{\C{recordedMutation.ts}\\\C{api.ts}}{集合の同一ID/本文の受付。\par GETと60秒write期限を分離。送信権とinstanceを確認しheaderへ固定。\par 自動POST再送ではない。}
\draw[call](single.west)--(ledger.east);\draw[call](single.east)--(mutation.west);
\File{guard}{appc}{.2}{4.9}{7.2}{\C{writeLease.ts}\\\C{serviceIdentity.ts}}{同originのWeb Locksで送信タブを一つにする。処理中でない担当へ明示的に引継ぎを依頼。\par HTTPと台帳処理の間は引き渡さず、編集/閲覧/3Dは保持。\par healthで台帳IDを確認し、別IDは自動採用しない。}
\File{service}{outc}{9.4}{4.9}{7.2}{\C{backend/app.py}\\\C{backend/application.py}}{headerと台帳instanceを照合。\par 受付とGET照会、\C{flight_execution}を公開。pool停止後の投入を止め、固定結果GETは保つ。\par 未完了は明示再開へ。}
\File{storage}{outc}{18.6}{4.9}{7.2}{\C{backend/storage.py}}{\C{service_identity}に永続UUID。\par 原spec・二段manifest・8成果物から完成物のDB登録漏れを回復。\par \C{run_recoveries}に旧state/errorと再計算なしを残す。}
\draw[call](single.south)--++(0,.35)-|(guard.north);
\draw[data](mutation.south)--++(0,.65)-|(service.north);
\draw[call](service.east)--(storage.west);
\Tip{meaning}{ink}{.2}{10.1}{25.6}{受付・実行・保存の意味を分ける}{応答不明は未受付の証明ではない。別instanceへ旧要求を移さない。同一台帳の複製や巻戻し、別ブラウザの排他制御までは保証しない。完成物回復の時刻は回復記録時刻である。操作はCOMMANDS C-09、HTTPと例外は実装ノートへ。}
\end{Map}

\Page{weatherc}{WEB-HISTORICAL}{原日時：関心から実場を選び、共通核へ渡す}{階層3 / 気象図 → 原日時入力 → historical計画 → ensembleの共通枝}{新しい飛行エンジンや別の結果画面は作らない。原日時の所有と場の同一性を明示し、既存の比較・履歴・干渉群へ接続する。}
\begin{Map}
\File{interest}{weatherc}{.2}{.3}{7.2}{\C{ClimateFlightHandoff.tsx}}{固定集計図の出典・期間/元半月ID・気象有効時刻・地域。\par 半月/月/季節の関心をAppへ。年間48区分は渡さない。\par 瞬時場や代表原日時を生成する機能ではない。}
\File{inputs}{appc}{9.4}{.3}{7.2}{\C{HistoricalEntry.tsx}\\\C{HistoricalPreparation.tsx}}{共通機体と原日時一覧を編集。\par 計画/台帳を確認して明示実行。\par \C{historicalDomain.ts}で入力・差分・支持を比較。\par 共通窓型は\C{historicalTypes.ts}。\C{ensembleDomain}も純型へ依存し、逆方向のimportを除く。}
\File{projection}{appc}{18.6}{.3}{7.2}{\C{historicalWorkspace.ts}}{共通projectの画面だけを用途別に絞る。\par Appの\C{season_real}状態。\par 保存から予報や旧結果を削らない。}
\draw[data](interest.east)--(inputs.west);\draw[data](inputs.east)--node[above,font=\scriptsize]{App経由}(projection.west);
\File{catalog}{weatherc}{.2}{5}{7.2}{\C{backend/historical_catalog.py}}{起動時path/hash/schema照合。\par ApplicationServiceのsourcesへ登録。原UTCと個々の場を保持。\par HTTP任意path入口ではない。}
\File{plan}{dync}{9.4}{5}{7.2}{\C{ensemble/historical_planning.py}}{原日時×共通機体caseを展開。\par 原UTCの重複と二重の時刻所有を拒否。未取得/支持不足は台帳へ。各trialに実場を固定。}
\File{existing}{outc}{18.6}{5}{7.2}{\Go{WEB-ENSEMBLE-BACKEND}{既存ensemble/の共通枝}}{\C{service.py} → 同じworker。\par \C{views.py} → 原日時台帳。\par 既存集計 → ForecastScreenの\C{ensembleProjection}とruntime。}
\draw[data](catalog.east)--(plan.west);\draw[data](inputs.south)--++(0,.55)-|node[pos=.25,above,font=\scriptsize,fill=white,inner sep=1pt]{API経由}(existing.north);\draw[call](existing.west)--(plan.east);
\Tip{set}{ink}{.2}{10}{25.6}{集合IDと一つの場のIDを混同しない}{\C{historical_window_set}は窓集合の指紋。支持境界と現在場との照合はmembers単位で行う。\C{parameter:null}を機体仮幅へ読み替えず、\C{weather_window}を原日時台帳へ返す。式・母集団はTHEORY、登録と操作はCOMMANDS、厳密な保存契約は実装ノートへ。}
\end{Map}

\Page{weatherc}{WEB-WEATHER}{weather：計画から完成場への公開を管理する}{階層3 / backend → weather → 契約・計画・索引・service}{取得順序と保存の寿命を扱うapplicationの枝である。物理式・GRIBの意味・気象場の補間はenvironmentへ委譲する。}
\begin{Map}
\File{parent}{weatherc}{.2}{.2}{25.6}{\C{backend/weather/}：固定要求 → 取得job → 完成asset}{読むもの：観測inventory、候補ごとの必要窓、領域、予算。渡すもの：保存計画と仕事の状態、検証済みsource。候補そのものや過去のrunは変更しない。}
\File{contract}{appc}{.2}{2.8}{7.2}{\C{contracts.py}}{InventoryRefresh / WeatherPlanRequest / AcquisitionRequest。\par 有限な型・範囲を受け入れる。\par 任意のURL/pathを受け取らない。}
\File{plan}{weatherc}{9.4}{2.8}{7.2}{\C{planning.py}}{make_plan / region_bounds。\par GeographicLibで矩形へ、GFS予定列で支持時刻へ。\par 依存・復号identity・メモリ見積。}
\File{catalog}{weatherc}{18.6}{2.8}{7.2}{\C{catalog.py}}{observe_inventory / Links。\par 固定一覧のrun/leadを有限に観測し、独立したday/runの失敗を記録。gatewayへ通信し、得た原HTMLをrecorderへ渡す。}
\draw[contain](parent.south)--++(0,.2)-|(contract.north);\draw[contain](parent.south)--(plan.north);\draw[contain](parent.south)--++(0,.2)-|(catalog.north);
\File{service}{outc}{.2}{6.7}{12}{\C{service.py / WeatherService}}{計画/仕事と一つの取得thread、取消/retry。\par \C{_raw_donors}：登録済み完成assetの一致rawを、\C{raw_provider}経由でgfsへ渡す。\par \C{_publish}：場・receipt・sourceを照合し、marker/DB登録後に完成場を親へ返す。}
\File{state}{outc}{14.1}{6.7}{11.7}{repo外：\C{state/weather/}}{registry.sqlite3：計画・job・asset参照。\par inventories/：原HTML・観測manifest。\par acquisitions/：raw・失敗・bundle・ASSET。\par marker単独でなく、DB登録後に利用可。}
\draw[imp](service.north)--(contract.south);\draw[call](service.north east)--++(0,-.5)-|(plan.south);\draw[call](service.north east)--++(0,-.5)-|(catalog.south);\draw[call](service.east)--node[above,font=\scriptsize]{保存/照合}(state.west);
\Tip{lifetime}{ink}{.2}{11.1}{25.6}{受付・途中・公開を分ける}{readyは取得の受付可。部分rawは診断/明示retry用に保持する。ASSETはpartialへ書き、rename後にhash照合とDB登録。完成assetは再利用しても旧runを変更しない。索引の観測・全飛行の支持・科学的精度は別の判断である。}
\end{Map}


\Page{weatherc}{WEB-CLIMATE-BACKEND}{climate：資料の違いを、固定結果へ結ぶ}{階層3 / backend → climate → service → 二つの資料adapter}{外生DBと元UTC標本を混ぜて一つの資料に見せない。資料別に同じsummary契約へ変換し、利用者が選んだ結果を別SQLiteへ固定する。}
\begin{Map}
\File{parent}{weatherc}{.2}{.2}{25.6}{\C{backend/climate/}：\C{ApplicationService.start}で起動 → 同期HTTP route}{sources / analyses POST / analysis GET。起動設定だけがDB/manifest pathを受ける。利用者queryはsource ID・範囲・気圧面・UTC時刻。取得用HTTPと飛行workerはここから呼ばない。}
\File{service}{outc}{.2}{3}{12}{\C{service.py / ClimateService}}{\C{sources / create_analysis / get_analysis}\par 複数資料をIDで選び、独立work lockで集計。\par 要求/内容の同一性を固定artifactへ結び、\C{climate.sqlite3}へ短いtransactionで保存。\par 保存GETは外生資料にも集計lockにも戻らない。}
\File{source}{weatherc}{14.1}{3}{11.7}{同じ入口、異なる資料の読取}{\C{source.py / ClimateSource}：提供DuckDB。\par \Go{WEB-CLIMATE-AGGREGATION}{保存平均とcountの結合}。任意の同支持原標本も扱う。\par \C{archive.py / WindArchiveSource}：\Go{WEB-CLIMATE-RAW}{独立した元UTC風標本}。\par 両者ともdescriptor/query/verify_fullを持つ。}
\draw[call](service.east)--(source.west);
\File{contract}{appc}{.2}{7.7}{7.2}{\C{contracts.py}}{\C{Bounds / ClimateQuery}\par 有限値・UTC集合・版識別。\par 全4時刻を省略形へ正規化。\par HTTP/SQL/保存を持たない。}
\File{samples}{weatherc}{9.4}{7.7}{7.2}{\C{samples.py}}{\C{MonthlySamples}\par manifest/配列/UTCと支持を照合。\par 完全月の加重経験分位を算出。\par 両adapterが呼ぶ。平均/分位から元風を合成しない。}
\Tip{lifetime}{outc}{18.6}{7.7}{7.2}{三つの寿命を分ける}{外生DB/配列：読取専用。\par adapter cache：有限・process内。\par SQLite artifact：固定・再訪可。\par 図の期間表示切替は再集計しない。}
\draw[imp](service.south west)--++(1,.35)-|(contract.north);\draw[call](source.south)--++(0,.35)-|(samples.north);
\Small{.2}{12.2}{26}{DuckDBは提供DB、NumPyは原配列と加重経験分位、SQLiteは固定結果保存。期間統計はenvironment.WeatherFieldを実装しない。相関を持つ飛行場や軌道MCへ直結したことにはならない。式はTHEORY、入力形式/例外/cacheは実装ノート、起動/取得はCOMMANDSへ。}
\end{Map}

\Page{weatherc}{WEB-CLIMATE-AGGREGATION}{climate：時間の結合と、地域の平均を分ける}{階層4 / source → aggregation・periods・extensions → statistics}{月や季節への変更は表示用の別母集団を明示する。平均量は件数から、風配は元の階級の件数から作り、架空の同時刻場や分位へ置き換えない。}
\begin{Map}
\File{host}{weatherc}{.2}{.2}{25.6}{\C{source.py}：資料の検査、支持の選択、結果の組立て}{DBとコードの指紋・変更を検査し、範囲/時刻の共通集計を有限cacheへ置く。選択面の地域値・風配を合わせ、\C{service.py}へ独立した結果を返す。図の粒度変更はこの処理を呼び直さない。}
\File{period}{weatherc}{.2}{3}{7.2}{\C{periods.py}}{\C{periods}：24半月 → 12か月 / 4季節。\par 元半月ID、UTC日数、対応時刻数。\par DJFは資料内の12/1/2月のpoolで、年をまたぐ連続冬列ではない。}
\File{aggregate}{weatherc}{9.4}{3}{7.2}{\C{aggregation.py}}{\C{cell_query / annual / spatial}\par 各格子のn重み時間平均 → Gaussian地域平均。件数/格子支持を検査。\par R/来向と格子平均の幅を再計算し、\C{scales}へ表示尺度を分ける。}
\File{extension}{weatherc}{18.6}{3}{7.2}{\C{extensions.py}}{\C{inspect_extensions / wind_rose}\par 追加tableの有無とschemaを検査。時刻別平均の支持面、単位/キー/件数を確認。\par 風配は固定1地点の階級countを結合。時刻subsetなら利用不可と返す。}
\draw[call](host.south)--++(0,.3)-|(period.north);\draw[call](host.south)--(aggregate.north);\draw[call](host.south)--++(0,.3)-|(extension.north);
\File{stats}{weatherc}{.2}{7.8}{12}{\C{statistics.py}}{\C{vector_summary / utc_half_months}\par 平均成分・速さ → R/FROM、UTC暦 → 元の件数。\par \C{combine_cells}は純粋な加重和によるSQLの照合先。\par aggregationとperiodsが使う。}
\Tip{support}{outc}{14.1}{7.8}{11.7}{支持と、保存結果の責務}{この提供DBでは全時刻38面、時刻別20面。欠測をゼロ補填しない。風配は地域外でも元地点を明示する。\par \C{summaries_by_grain}に月/季節を保存。旧artifactに新集計を後付けしない。}
\draw[imp](aggregate.south)--++(0,.4)-|(stats.north);\draw[imp](period.south)--++(0,.4)-|(stats.north west);
\Small{.2}{12.15}{26}{構造の責務はこの頁、n重み・Gaussian・経験頻度とUTC/JSTの意味はTHEORYへ。この枝は提供平均と風配countを結合する。月/地域分位はここで平均せず、同支持の原標本がある場合だけ\Go{WEB-CLIMATE-RAW}{samples}へ渡す。}
\end{Map}

\Page{weatherc}{WEB-CLIMATE-RAW}{原標本：取得、検査、分位の担当を分ける}{階層4 / tools → 外生manifest → archive → samples・periods・statistics}{元時刻u/vの束が初めて、任意の選択格子とUTC時刻を同じ分布へまとめる根拠になる。取得プログラムはサービスの起動・固定結果保存・画面を所有しない。}
\begin{Map}
\File{acquire}{toolc}{.2}{.3}{7.2}{\C{tools/}：製品別取得}{\path{acquire_jra3q_climate_samples.py}\par\path{acquire_era5_climate_samples.py}\par 完了月と出典を保存。\C{prepare_climate_demo.py}は元格子の小標本化/固定束の展開を通信なしで担当。}
\File{bundle}{weatherc}{9.4}{.3}{7.2}{repo外：原標本束}{\C{balloon.wind-samples/1}\par provider/年/格子/面/重み。\par 月ごとのUTC列とfloat32 u/v配列、bytes/SHA、出典。\par JRAとERAを同じ格子へ補間しない。}
\File{archive}{weatherc}{18.6}{.3}{7.2}{\C{archive.py / WindArchiveSource}}{全12か月の原標本adapter。\par \C{_moments / _annual}は全高度の年間値を4母集団cacheへ。\par \C{_spatial}は選択面、\C{_rose_counts}は1点を一巡し3粒度へ。}
\draw[data](acquire.east)--(bundle.west);\draw[data](bundle.east)--(archive.west);
\File{samples}{weatherc}{.2}{5.2}{12}{\C{samples.py / MonthlySamples}}{\C{assert_current / verify_full / profiles}\par 起動時に完全UTC列・shape・有限値・SHAを照合。\par 1面ずつu/vを読み、速さをNumPy加重分位へ。\par 前後半medianと月p10/p90は別々の母集団。\par 旧DBに結び付けたraw-months/1も同じ実体で検査する。}
\File{common}{weatherc}{14.1}{5.2}{11.7}{共通の集計の子}{\C{periods.py}：UTC暦・件数。年間48区分は原標本から直接計算し、旧半月/月/季節を保つ。\par \C{statistics.py}：格子重み平均・R/FROM。\par \C{aggregation.scales}：表示尺度。\par \C{archive.isa_height}：表示専用ISA参考軸。\par 幾何高度未取得を実高度0へ置き換えない。}
\draw[call](archive.south)--++(0,.35)-|(samples.north);\draw[imp](archive.south)--++(0,.35)-|(common.north);
\Tip{same}{outc}{.2}{10.1}{12}{同じ結果形式でも、母集団は別}{提供DBは2016--25年。新原標本例は2024年。\par 全域JRAは38面/570格子、ERA5は37面/1188格子。配布用2×2元格子例は機能確認用で、全域の代表標本ではない。\par source ID、年、native座標、重み、出典を保つ。}
\Tip{boundary}{ink}{14.1}{10.1}{11.7}{この枝の外にある仕事}{取得進捗と実測時間はS36。数式はTHEORY。\par 取得CLI/再開はCOMMANDS。\par 月平均や分位から瞬時共同場を作る処理、年の自由選択、飛行MCはこの枝にない。}
\end{Map}

\Page{weatherc}{WEATHER-FILE}{environment：製品と場の意味を分ける}{階層2 / balloon_sim → environment → 場・製品・保存・通信}{親の責務は「保存入力を、必要な時空間支持を持つ環境として渡すこと」。圧力面/モデル面の再構成と、製品固有の時刻・座標契約を別の実体にする。}
\begin{Map}
\File{parent}{weatherc}{.2}{.2}{25.6}{\C{environment/}}{読むもの：取得要求または保存bundle。渡すもの：場、または原本/来歴付き保存物。利用元：CLI・backendの取得service/飛行worker・fixture生成。}
\File{fields}{weatherc}{.2}{3}{7.2}{\Go{WEATHER-FIELD}{\C{fields.py}}}{\C{PressureLevelField}\par 共通高度の照会、支持/欠測、高度変換。製品時刻規則は持たない。}
\File{bundle}{weatherc}{9.4}{3}{7.2}{\C{bundle.py}}{\C{WeatherField}\par 既存schemaを受ける入口。場構築にGFS時刻規則を追加する。}
\File{rules}{weatherc}{18.6}{3}{7.2}{\C{gfs_contract.py}}{33圧力面、lead、配布URL、\C{validate_gfs_times}\par GFSの固定run/予定時刻を検査。}
\File{storage}{weatherc}{.2}{6.6}{7.2}{\C{storage.py}}{\C{read_json / write_bundle}\par \C{load_weather}はschemaでGFS / JRA strict / JRA暫定Bを選ぶ。strict JSON/gzip。}
\File{gfs}{weatherc}{9.4}{6.6}{7.2}{\Go{WEATHER-ACQUIRE}{\C{gfs.py}}}{取得・raw照合・復号。\par 製品仕様をGFS規則へ、bundle検証を場入口へ委譲。}
\File{network}{weatherc}{18.6}{6.6}{7.2}{\C{nomads.py / NomadsGateway}}{一覧/GRIBの固定origin通信。\par CLI/APIで共有するOS-user lock、間隔・取消・時間/bytes制限。\par 気象の復号規則は持たない。}
\draw[contain](parent.south)--++(0,.3)-|(fields.north);\draw[contain](parent.south)--(bundle.north);\draw[contain](parent.south)--++(0,.3)-|(rules.north);
\draw[contain](parent.west)--++(-.18,0)|-(storage.west);\draw[contain](parent.east)--++(.18,0)|-(network.east);\draw[contain](parent.south)--++(0,.3)-|(8.85,6.2)-|(gfs.north);
\draw[imp](bundle.west)--(fields.east);\draw[imp](bundle.east)--(rules.west);
\draw[imp](gfs.east)--(network.west);
\File{model}{weatherc}{.2}{10.3}{12}{モデル面の再構成}{\Go{WEATHER-MODEL}{\C{model_levels.py}}：元列ごとのstrict照会。\par \Go{WEATHER-SURFACE}{\C{model_surface.py}}：原UTC別の仮想列と地上橋。}
\File{jra}{weatherc}{14.1}{10.3}{11.7}{\C{jra3q.py}：製品と明示方針}{\Go{WEATHER-MODEL}{\C{Jra3qModelField}}：strict。\par \Go{WEATHER-SURFACE}{\C{Jra3qSurfaceField}}：暫定B。\par 同じ原UTC・係数・full pを束縛。}
\draw[imp](jra.west)--node[above,font=\scriptsize]{継承・構築}(model.east);
\draw[contain](parent.west)--++(-.25,0)|-(model.west);\draw[contain](parent.east)--++(.25,0)|-(jra.east);
\end{Map}

\Page{weatherc}{WEATHER-FIELD}{圧力面場：構築してから照会する}{階層3 / environment → fields.py / bundle.py → クラスとメソッド}{入力の形を受け入れる段階と、一つの時空間点を評価する段階を分ける。値の意味は理論ガイドの気象再構成節へ戻る。}
\begin{Map}
\File{bundle}{weatherc}{.2}{.2}{7.2}{\C{bundle.WeatherField}}{\C{PressureLevelField}を継承。\par 元のGFS時刻列を追加検査する。}
\File{base}{weatherc}{10}{.2}{15.8}{\C{fields.PressureLevelField}}{保持：times / latitudes / longitudes / pressures / fields / surface / metadata\par 共通の支持・高度列・必要量を扱い、原配布先を知らない。}
\draw[imp](bundle.east)--node[above,font=\scriptsize]{継承}(base.west);
\Card{init}{weatherc}{.2}{3.6}{7.2}{\C{__init__}}{軸・配列形・有限値・高さ列を検査。\par tuple/deepcopyで元辞書から分離。}
\Card{sample}{weatherc}{9.4}{3.6}{7.2}{\C{sample}}{要求量→支持列→各列の鉛直再構成→水平/時間合成。\par 返値：量・地形・品質。}
\Card{ground}{weatherc}{18.6}{3.6}{7.2}{\C{ground_altitude}}{同じ支持の地形Hを混合し、幾何高度mへ返す。}
\draw[contain](base.south)--++(0,.45)-|(init.north);\draw[contain](base.south)--++(0,.45)-|(sample.north);\draw[contain](base.south)--++(0,.45)-|(ground.north);
\Card{support}{weatherc}{.2}{7.1}{12}{\C{_support / _mix}}{前者が非ゼロの(time,lat,lon,weight)、後者が地上配列の重み付き値を作る。sampleと地形照会が共用する。}
\Card{helpers}{weatherc}{14.1}{7.1}{11.7}{共有ヘルパー}{\C{_axis / _array / _bracket}\par \C{_number / _utc / _required_object}\par 2つの高度変換、WeatherError。}
\hypertarget{WEATHER-SAMPLE}{}
\Tip{meaning}{weatherc}{.2}{10.6}{25.6}{sampleの中で行うこと、飛行へ渡すこと}{列ごとに地下面を除き、地上アンカーまたは有効圧力面から共通高度へ再構成する。その実値を混ぜる。flag/checked/upperはこのメソッドに閉じる小補助。寄与列の地表延長を品質として返し、飛行の相・接地イベントは扱わない。}
\end{Map}

\Page{weatherc}{WEATHER-MODEL}{strictモデル面場：元列の支持を守る}{階層3 / environment → model_levels.py・jra3q.py → クラスとメソッド}{飛行核のsample/ground_altitudeを共用する。GFSへ型を偽装せず、JRAの鉛直座標と厳密支持を明示した枝である。暫定Bでも、この製品検査と原列の保持を土台に使う。}
\begin{Map}
\File{jra}{weatherc}{.2}{.2}{12}{\C{jra3q.Jra3qModelField}}{\C{__init__}：schema/product、6h、100層、既知係数を検査。\par \C{full_pressure_pa}でSPと係数から各列pを構成。\par 元時刻・出典・境界方針をmetadataに残す。}
\File{model}{weatherc}{14.1}{.2}{11.7}{\C{model_levels.ModelLevelField}}{保持：times / latitudes / longitudes / levels\par heights / pressures / fields / surface / metadata\par tuple/deepcopyで入力から分離。製品を知らない。}
\draw[imp](jra.east)--node[above,font=\scriptsize]{継承}(model.west);
\Card{sample}{weatherc}{.2}{4.1}{7.2}{\C{sample}}{要求量 → 各列H支持 → 鉛直補間 → 水平/時間合成。\par field値、品質、補間地形を返す。}
\Card{support}{weatherc}{9.4}{4.1}{7.2}{\C{_support / _ground}}{UTC/水平の正重みを確定。\par 同じ支持で地形Hを合成。\par 同値経度±360のみ許容。}
\Card{ground}{weatherc}{18.6}{4.1}{7.2}{\C{ground_altitude}}{地形の共通API。\par 合成Hを幾何ASL mへ変換。\par 地形を最下層へ置換しない。}
\draw[contain](model.south)--++(0,.5)-|(sample.north);\draw[contain](model.south)--++(0,.5)-|(support.north);\draw[contain](model.south)--++(0,.5)-|(ground.north);
\File{helpers}{weatherc}{.2}{8}{12}{共有依存：\C{fields.py}}{軸・配列・値・UTC検査、bracket、高度変換、WeatherError。\par GFSの\C{PressureLevelField.sample}へ委譲しない。}
\File{prepare}{toolc}{14.1}{8}{11.7}{固定原本の変換：\C{tools/prepare_jra3q_flight_fixture.py}}{旧検証loader → 新bundle → 保存/全payload読戻し。\par runtimeはこのtoolをimportしない。一般取得器は後続。}
\draw[imp](support.south)--++(0,.5)-|(helpers.north);
\Small{.2}{11.6}{26}{呼出順：CLI → storageのschema dispatch → 同じflight核 → 結果保存。保存B例のGUI入口は\Go{WEB-SERVICE}{共通service}へ。任意過去窓の取得画面は未実装。strictでは元列の全正重み支持を要求し、地表不足は停止結果になる。地上解析を明示的に使う別枝は\Go{WEATHER-SURFACE}{次頁の暫定B}。数式はTHEORY、旧strictの操作はCOMMANDS C-10へ。}
\end{Map}

\Page{weatherc}{WEATHER-SURFACE}{暫定B：原UTCの仮想列から地上へつなぐ}{階層3 / environment → jra3q.py → model_surface.py → 保持するstrict場}{旧strictのsampleを書き換えず、別schemaとpolicyで選ぶ。通信・取得UIはこの枝に置かない。Aの地形列延長は外側の比較実験に残す。}
\begin{Map}
\File{entry}{weatherc}{.2}{.2}{12}{\C{jra3q.Jra3qSurfaceField}}{\C{balloon.weather.jra3q_surface/1}とproductを検査。\par \C{reconstruction.policy / join_model_levels}を必須化。\par 原配列で\C{Jra3qModelField}を構築し、Bへ渡す。}
\File{base}{weatherc}{14.1}{.2}{11.7}{\C{model_surface.ModelSurfaceField}}{Bの照会を実装。\C{source}に検査済みstrict場を保持。\par 4地上解析配列と変数別固定joinを受ける。\par 入力配列を別再構成で上書きしない。}
\draw[imp](entry.east)--node[above,font=\scriptsize]{継承}(base.west);
\Card{construct}{weatherc}{.2}{4.2}{7.2}{\C{__init__ / _validate_fixed_joins}}{全原UTC・全列の地形/接続高度を検査。\par p/T/qはlevel1、u/vはlevel2。\par pは地表、T/qは2m、風は10mの節点より上でなければ拒否。}
\Card{sample}{weatherc}{9.4}{4.2}{7.2}{\C{sample / _column / _mix}}{原UTCでnative H/Fを水平合成。\par 同じ絶対Hへ鉛直照会。\par 最後にUTC間を線形混合。\par 必要値を平均前にも検査する。}
\Card{ground}{weatherc}{18.6}{4.2}{7.2}{\C{ground_altitude / metadata}}{地形APIは保持するsourceへ委譲。\par policy・不使用のnative1風・上空にも及ぶ変更を記録。\par 地下の有限延長は行わない。}
\draw[contain](base.south)--++(0,.5)-|(construct.north);\draw[contain](base.south)--++(0,.5)-|(sample.north);\draw[contain](base.south)--++(0,.5)-|(ground.north);
\File{source}{weatherc}{.2}{8.7}{12}{包含する原列場：\Go{WEATHER-MODEL}{\C{Jra3qModelField}}}{元UTC・全native面・列別full p・地形・来歴を保持。\par \C{_support / _ground / _check_value}をBから共用する。\par 元列の\C{sample}へBの値計算を委譲するのではない。}
\Tip{contract}{weatherc}{14.1}{8.7}{11.7}{同じ飛行protocol、異なる場の意味}{\C{sample}は必要量・品質・地形・橋の幅を返す。\par \C{storage.load_weather} → この場 → 既存simulate。\par 最低風の置換だけでなく、上空の値/上端支持も変わり得る。精度・保存性の保証はしない。}
\draw[imp](sample.south)--++(0,.4)-|(source.north);\draw[imp](ground.south)--++(0,.6)-|(source.east);
\Small{.2}{12.5}{26}{方針の正本は\C{native_horizontal_first_surface_v1}。2m/10m以下の定値、B1圧力、要求値欠測と仮想列の支持はTHEORYへ。保存schemaの自動置換・GFSへのfallback・過去場の取得画面は実装していない。}
\end{Map}

\Page{weatherc}{WEATHER-ACQUIRE}{GFS：rawを確定して共通の復号へ渡す}{階層3 / environment → gfs.py → acquire / replay / decode}{新規取得・明示resumeと原本再生は、rawの用意と照合で分かれる。復号する量・格子・時刻の意味は同じdecodeで受け入れる。}
\begin{Map}
\Card{acquire}{weatherc}{.2}{.3}{10.5}{\C{acquire_gfs(request, output_dir)}}{\C{resume / progress / cancel / raw_provider}。\par 同jobのprefix → 提供候補の照合/独立copy → 不足rawの通信。\par 全窓をdecodeへ。backendのDBを直接読まない。}
\Card{replay}{weatherc}{15.3}{.3}{10.5}{\C{replay_gfs(acquisition_dir, output_path)}}{保存要求・lead・URL・名前・bytes・SHAを照合。\par 通信で不足を補わず、同じrawを再復号。}
\Card{decode}{weatherc}{7.7}{4.2}{11}{\C{decode_gfs(paths, specification, provenance)}}{ecCodes/NumPyを遅延import。\par 製品・run・valid・格子・単位・風基準・必要量を確認。\par 配列を整列し、WeatherFieldでbundleを検査。}
\draw[call](acquire.south)--(decode.north);\draw[call](replay.south)--(decode.north);
\File{network}{weatherc}{.2}{8.4}{7.2}{\C{NomadsGateway}}{\C{fetch / download_to}\par 前応答後10秒以上、Retry-After、時間/bytes制限。\par OS-user lockでCLI/APIを共通制御。}
\File{write}{weatherc}{9.4}{8.4}{7.2}{\C{storage.write_bundle}}{返ったbundleを入口側が渡す。\par 決定的JSON/gzip、新規出力のみ。}
\File{field}{weatherc}{18.6}{8.4}{7.2}{\C{bundle.WeatherField}}{正規化した結果の受入。\par decodeが試算飛行を行うことはない。}
\draw[data](decode.south)--(write.north);\draw[imp](decode.east)--++(1,0)-|(field.north);
\draw[call](acquire.west)--++(-.15,0)|-node[pos=.8,above,font=\scriptsize]{raw通信}([yshift=5mm]network.north)--(network.north);
\Small{.2}{12.2}{26}{\C{_verified_raw_candidates / _copy_verified_raw}はrawの内容とコピー先を照合。候補の登録責務は\Go{WEB-WEATHER}{weather/service}側に残る。通信は\C{_download → NomadsGateway}、時刻/層/URLはgfs_contract、復号見積はestimate_decode_bytes。操作/障害時はCOMMANDS・実装ノートへ。}
\end{Map}

\Page{dync}{ENSEMBLE-CORE}{ensemble：元標本と共通集計を、入出力から分ける}{階層2 / balloon_sim → ensemble → sampling・statistics}{この枝は純粋な値の変換を担う。抽選した元値を完全設定へ展開し、固定結果を同じ定義で集計する。仕事の投入、DB、HTTP、ファイルの公開はbackend側にある。}
\begin{Map}
\File{package}{dync}{.2}{.2}{25.6}{\C{balloon_sim/ensemble/}：\C{__init__.py}が公開名をまとめる}{入口は\C{freeze_drawset / resolve_config / analyze_results}。一飛行の\C{simulate}を内蔵せず、予定標本・個別飛行・表示集計の寿命を分ける。}
\File{sampling}{dync}{.2}{3}{12}{\C{sampling.py}}{\C{freeze_drawset}：明示した元量・単位・一様幅・理由・seedから、元実現値とdraw IDを固定。\par \C{resolve_config}：選んだ物理元量を一つのconfigへ展開し、全設定を再検証。}
\File{statistics}{outc}{14.1}{3}{11.7}{\Go{ENSEMBLE-ANALYSIS}{\C{statistics.py}}}{\C{analyze_results}：全台帳のIDと選択を検査し、選択群を着地幾何と相別履歴へ渡す。\par null結果や履歴なし停止も、選択母数から消さない。}
\draw[contain](package.south)--++(0,.3)-|(sampling.north);\draw[contain](package.south)--++(0,.3)-|(statistics.north);
\File{config}{dync}{.2}{7.4}{7.2}{既存\C{flight/config.py}}{\C{validate_config / FlightError}\par 展開先の型・モデル・値域を担う。\par モデル不活性を入力欄の追加で隠さない。}
\File{rng}{toolc}{9.4}{7.4}{7.2}{外部：\C{NumPy}}{\C{Generator / PCG64 / uniform}\par 乱数列生成は既存ライブラリへ。\par seedだけでなく実現値・版を保存側へ返す。}
\Card{outputs}{outc}{18.6}{7.4}{7.2}{固定する値 → \Go{WEB-ENSEMBLE-BACKEND}{backend}}{drawsetと完全config → 計画。\par 完了結果/選択ID → 共通分析。\par 気象の取得や飛行の再実行は、この枝から起こさない。}
\draw[imp](sampling.south west)--++(1,-.1)|-(config.north);\draw[imp](sampling.south east)--++(-1,-.1)|-(rng.north);\draw[data](statistics.south)--++(0,.5)-|(outputs.north);
\Small{.2}{12.1}{26}{現在の元量は等温ガス質量・高度破裂・基準降下速度。搭載質量の片側だけの抽選や製品の分布推定は実装していない。方式と式はTHEORY、受付予算とAPIは実装ノートへ。}
\end{Map}

\Page{outc}{ENSEMBLE-ANALYSIS}{集合の集計：幾何と時刻の二つの枝へ降りる}{階層3 / ensemble → statistics → landing・histories・geography}{読む対象は同じ固定結果と選択IDである。母数、着地点、相付き記録の違いを、図の描画から独立した共通出力へ残す。}
\begin{Map}
\File{entry}{outc}{.2}{.2}{25.6}{\C{statistics.py / analyze_results}}{選択ID・科学結果の外枠を照合 → 選択行を二つの解析へ。選択数、結果あり、履歴あり、欠落を別に返す。run/attempt/sourceのhash照合は呼出側が担当する。}
\File{landing}{outc}{.2}{3}{12}{\C{landing.py / landing_analysis}}{着地のみ → 局所平面 → 共分散/固有軸 → 経験50/90/95\%と全点extent。\par 点/線/面と非支持理由を返す。含有IDを保持し、地図の輪郭交差から干渉数を作らない。}
\File{histories}{outc}{14.1}{3}{11.7}{\C{histories.py / history_analysis}}{\C{_metric_records}で原記録を量と相へ分ける。\par 同じ経過時刻で相内だけを内挿し、平均/10--90\%帯と各量の有効Nを返す。原recordsは変更しない。}
\draw[call](entry.south)--++(0,.3)-|(landing.north);\draw[call](entry.south)--++(0,.3)-|(histories.north);
\File{geo}{weatherc}{.2}{7.4}{7.2}{\C{geography.py / LocalPlane}}{\C{xy / lonlat / supported_ring}\par WGS84局所方位距離座標と地理座標を相互変換。\par GeographicLibを使用。}
\File{libs}{toolc}{9.4}{7.4}{7.2}{外部の一般計算}{landing：\C{numpy.linalg.eigh}\par \C{scipy.spatial.ConvexHull}\par histories：\C{numpy.interp}\par \C{quantile / mean}。}
\Tip{phase}{dync}{18.6}{7.4}{7.2}{相と終端の意味}{上昇左極限/下降右極限を相別に保持。終端の後へ外挿しない。\par 同時刻二相のNは足さない。\par 平均位置線は実在の一飛行ではない。}
\draw[imp](landing.south west)--++(1,0)|-(geo.north);\draw[imp](landing.south east)--++(-1,0)|-(libs.north);\draw[imp](histories.south west)--++(1,0)-|(libs.north east);
\Small{.2}{12.1}{26}{局所幾何の非支持は着地そのものの失敗ではない。禁止領域への全着地点判定はfrontend/math、選択ID・領域版・分析artifactの保存はbackendへ戻る。詳しい経験域/内挿の定義と許容差はTHEORYに置く。}
\end{Map}

\Page{dync}{DYNAMICS-FILE}{flight：一飛行の意味を束ねる}{階層2 / balloon_sim → flight → 設定・物理・イベント・数値・進行}{親の入力はconfigと場、出力は一飛行の結果辞書。ここでは取得元やファイル形式を知らず、飛行中の判断に専念する。}
\begin{Map}
\File{trajectory}{dync}{9.4}{.2}{7.2}{\Go{DYNAMICS-SIMULATE}{\C{trajectory.py}}}{\C{simulate(config, weather)}\par 相・場照会・イベント優先・採用記録を組み立てる。}
\File{config}{dync}{.2}{3.8}{7.2}{\C{config.py}}{\C{validate_config / FlightError}\par 型・単位・モデル組合せ・既定値。\par 検証した設定コピーを返す。}
\File{models}{dync}{9.4}{3.8}{7.2}{\Go{FLIGHT-MODELS}{\C{models.py}}}{\C{compile_models / vertical_state}\par 気体・密度・浮力抗力と必要量を束縛。\par 取得も積分も行わない。}
\File{events}{dync}{18.6}{3.8}{7.2}{\Go{FLIGHT-MODELS}{\C{events.py}}}{\C{BurstEvent}\par 高度/径の条件、診断量、既知高度への正確な到達刻み。\par 相の切替はtrajectory側。}
\File{numerics}{dync}{9.4}{7.6}{7.2}{\Go{FLIGHT-NUMERICS}{\C{numerics.py}}}{右辺と状態の試行、符号交差の探索。\par 気象製品・気球の式は知らない。}
\draw[imp](trajectory.south)--++(0,.5)-|(config.north);\draw[imp](trajectory.south)--(models.north);\draw[imp](models.east)--(events.west);\draw[imp](trajectory.east)--++(.8,0)|-(numerics.east);
\Tip{boundary}{dync}{.2}{11}{25.6}{モデル追加と数値変更は別の判断}{Re依存は抗力則、慣性や熱は状態と方程式の変更。今回は現物の責務を分離し、任意モデル登録・任意次元の飛行モデルを完成したとはしない。将来の構成はENV、現在選べるmodeは実装ノートとCOMMANDSで確認する。}
\end{Map}

\Page{dync}{DYNAMICS-SIMULATE}{trajectory：場とモデルを一飛行へ接続する}{階層3 / flight → trajectory.py → simulate と一飛行の局所状態}{時刻・相・採用記録は一飛行に属する。数値算術と物理の評価を外へ委譲したうえで、これらの進行を一箇所で判断する。}
\begin{Map}
\Card{driver}{dync}{.2}{.2}{25.6}{\C{simulate}の進行}{設定受入・compile_models → 放球状態 → 試行 → 破裂/地面の先着を選択 → 採用記録 → 相切替または終端。支持を失った場合は、採用済みの結果と停止理由を保つ。}
\Card{query}{weatherc}{.2}{3.2}{7.2}{場照会の局所補助}{\C{stamp / ground / sample_at}\par 経過秒をUTCへ。必要量をモデルから取得。試行点の地表延長を区別する。}
\Card{rhs}{dync}{9.4}{3.2}{7.2}{右辺と試行}{\C{rhs}：球面移流＋鉛直式\par \C{AdaptiveRK45.trial}へ委譲\par 共通の場を積分段ごとに照会。}
\Card{event}{dync}{18.6}{3.2}{7.2}{イベントへの接続}{\C{burst_value / locate_crossing}\par 束縛済みBurstEventへ条件評価、numericsへ根探索を渡す。地面条件は局所ground。}
\draw[call](rhs.west)--(query.east);\draw[call](event.south)--++(0,.35)-|node[pos=.25,below,font=\scriptsize]{径破裂で必要な場を照会}(query.south);
\Card{record}{outc}{.2}{7.2}{12}{採用点と破裂記録}{\C{make_record / accept / make_burst_event}\par make_recordの検証後にacceptが時刻・状態・記録を一組で採用。到達した破裂点を保ち、下降開始を別に受け入れる。同時刻2相を区別。}
\Card{failure}{dync}{14.1}{7.2}{11.7}{失敗と共有状態}{\C{failure}がcode・時刻を持つ停止理由へ。\par phase / elapsed / state / records / events はsimulateの局所状態で、別飛行とは共有しない。}
\Tip{exit}{dync}{.2}{10.8}{25.6}{親へ返す結論}{降下中の地面交差がlanded。上昇衝突・支持喪失・上限超過はstopped。結果辞書を返すまでが責務で、ファイル保存はresults。数値法を替える際も、相・停止・採用点の意味を無言で変えない。}
\end{Map}

\Page{dync}{FLIGHT-MODELS}{物理とイベント：値を評価し、進行は返す}{階層3 / flight → models.py / events.py → 現在の式と必要量}{上昇・降下の法則、破裂条件、相を進める操作を分ける。式の導出と仮定は理論ガイドに置き、ここではそれを担う実体を示す。}
\begin{Map}
\File{model}{dync}{.2}{.2}{12}{\C{models.py}}{読む：検証したconfig。compile_modelsが各相の入力を分離。\par 渡す：FlightModels(ascent,descent,burst)。各相は評価関数と必要量を持つ。}
\File{events}{dync}{14.1}{.2}{11.7}{\C{events.py}}{読む：独立した閾値・gas_state callable・指定上昇速度。\par 渡す：破裂余裕・診断・正確な刻み・高度丸め。相の設定配置を知らない。}
\Card{vertical}{dync}{.2}{3.8}{7.2}{\C{vertical_state}}{上昇：規定速度 / 等温浮力\par 降下：CdA / 基準速度\par 密度を使う経路だけ密度を評価。}
\Card{gas}{dync}{9.4}{3.8}{7.2}{\C{_gas_state / _density}}{前者：体積・直径・投影面積\par 後者：\C{moist_air_density}へ\par 選んだ仮定を共通に評価する。}
\Card{burst}{dync}{18.6}{3.8}{7.2}{\C{BurstEvent}}{value / diagnostics / exact_step\par project_altitude、必要量の有無。\par gas_stateを引数で受ける。modelsへ逆importしない。}
\draw[call](vertical.east)--(gas.west);\draw[call](burst.south)--++(0,.4)-|node[pos=.25,below,font=\scriptsize]{注入した気体計算のcallable}(gas.south);
\Card{needs}{weatherc}{.2}{7.5}{12}{\C{PhaseModel / FlightModels}}{compile_modelsが各相のevaluateとrequired_fieldsを束縛する。trajectoryはmode名で式を選ばず、phaseと評価契約だけを使う。熱等の状態追加は別の契約拡張が必要。}
\Tip{change}{dync}{14.1}{7.5}{11.7}{変更時に一緒に見る境界}{法則の変更 → configと必要量\par 状態の追加 → trajectory/結果schema\par 径の式変更 → 破裂と診断出力\par 同じ計算自由度を二重に独立入力しない。}
\Small{.2}{12}{26}{現役の単純経路を残して比較する。熱・慣性・Re依存をまだ選べるようにはしていない。等温仮定の採用条件は後から拡張できる設計であり、現時点の感度試験を再要求しない。}
\end{Map}

\Page{dync}{FLIGHT-NUMERICS}{numerics：一般数値法をSciPyへ委譲する}{階層3 / flight → numerics.py → SciPy RK45 と brentq}{物理・場・相の意味は飛行側が持つ。一般的な誤差推定、適応刻み、区間補間、根探索は既存ライブラリの公開APIへ渡す。}
\begin{Map}
\File{caller}{dync}{.2}{.3}{7.2}{\C{trajectory.py}}{rhs・許容誤差・刻み上限を渡す。\par 支持の拒否、相切替、イベント優先、採用記録は呼出側。}
\Card{step}{dync}{9.4}{.3}{7.2}{\C{AdaptiveRK45}}{trial → \C{IntegrationStep}\par reset：破裂/拒否後の再初期化。\par state_atで区間内の状態を照会。}
\Card{root}{dync}{18.6}{.3}{7.2}{\C{locate_crossing}}{区間とイベント関数から時間根へ。\par \C{event_tolerance_s}は時間xtol。\par \C{numerical_metadata}が方法/版/設定を来歴へ返す。}
\draw[call](caller.east)--(step.west);\draw[data](step.east)--node[above,font=\scriptsize]{区間}(root.west);
\File{rk}{toolc}{.2}{4.3}{12}{外部：\C{scipy.integrate.RK45}}{step：局所誤差で刻みを調節。dense_output：採用区間の補間。\par 独自のRK4は現行本体から外し、比較元に保持。}
\File{brent}{toolc}{14.1}{4.3}{11.7}{外部：\C{scipy.optimize.brentq}}{既に挟んだ連続なイベント関数の根を求める。\par 地形・破裂の選択、停止結果の意味は持たない。}
\draw[imp](step.south)--(rk.north);\draw[imp](root.south)--(brent.north);
\Tip{contract}{dync}{.2}{8.4}{12}{飛行側へ戻す境界}{採用区間の時刻・状態・dense補間を返す。失敗試行を採用点にしない。非有限/solver失敗はIntegrationError。採用点は不等間隔で、固定出力時刻への再標本化は未実装。}
\Tip{scale}{toolc}{14.1}{8.4}{11.7}{許容誤差の単位を忘れない}{rtolと、緯度/経度degree・高度mのatol。\par max_stepは間隔上限。根精度や局所誤差の制御を着地点の精度保証へ拡張しない。独立試験はtest_flight_numerics.pyへ。}
\end{Map}

\Page{outc}{CLI-EXPORT}{results：意味を変えず、読める結果にする}{階層2 / balloon_sim → results → export.py / report.py}{読むものは完了・途中停止を含む結果と入力来歴。ここで物理を再計算したり、停止を成功へ変えたりしない。}
\begin{Map}
\File{export}{outc}{.2}{.3}{12}{\C{export.py}}{\C{export_result / make_geojson}\par 入力控え、結果、CSV、GeoJSON、来歴、manifestを新しい出力先へ保存。}
\File{report}{outc}{14.1}{.3}{11.7}{\C{report.py}}{\C{render_report}\par 結果をHTML文字列へ。軌道/高度SVG、時刻スライダー、設定、来歴。}
\draw[call](export.east)--(report.west);
\Card{geo}{outc}{.2}{4.4}{7.2}{\C{make_geojson}}{\C{_point}で経度・緯度・高度の順に変換。線とイベント点を返す。自分では保存しない。}
\Card{files}{outc}{9.4}{4.4}{7.2}{保存と指紋}{\C{sha256 / _json / _write_json}\par 入力と再帰した本体全pyの指紋、数値法/依存版。\par manifest自身はhash対象外。}
\Card{svg}{outc}{18.6}{4.4}{7.2}{表示の小補助}{\C{_trajectory_svg}\par \C{_altitude_svg}\par 後者のxyは図上座標変換。計算の座標式とは別。}
\draw[call](export.south)--(geo.north);\draw[call](export.south)--(files.north);\draw[call](report.south)--(svg.north);
\Tip{purpose}{outc}{.2}{8.8}{25.6}{表示を変更するときの戻り先}{この旧単一HTMLの水平図は背景地図なしの局所投影である。現在の比較地図はscreens/forecastの別表示。集合の経験域はensembleの共通分析から作り、実世界の事故確率とは区別する。単一HTMLの表示を、集合の図と同一の責務にしない。形式とhashの検査、画面の視認、科学的精度は別に評価する。}
\end{Map}

\Page{appc}{PROGRAM-CLI}{CLI：利用者の操作を接続する}{階層2 / balloon_sim → cli.py → main}{\C{python -m balloon_sim}から一つの操作を行う。公開していた整形/保存名は互換再exportで維持し、本体定義はresultsへ移した。}
\begin{Map}
\hypertarget{CLI-MAIN}{}
\File{main}{appc}{.2}{.3}{25.6}{\C{main(argv=None)}}{argparseで4操作を選択し、入力と既存出力を確認する。呼出結果をstdoutへ、エラーをstderrへ、終了コードを呼出元へ返す。}
\Card{fetch}{weatherc}{.2}{3.5}{5.5}{fetch}{read_json → acquire_gfs\par 通信と復号を行う。}
\Card{replay}{weatherc}{7.3}{3.5}{5.5}{replay}{replay_gfs\par 保存rawを照合し再復号。}
\Card{inspect}{weatherc}{14.4}{3.5}{5.2}{inspect}{load_weather\par 場のmetadataを表示。}
\Card{simulate}{dync}{21.2}{3.5}{4.6}{simulate}{load_weather + 設定\par simulate → export_result}
\draw[call](main.south)--++(0,.5)-|(fetch.north);\draw[call](main.south)--++(0,.5)-|(replay.north);\draw[call](main.south)--++(0,.5)-|(inspect.north);\draw[call](main.south)--++(0,.5)-|(simulate.north);
\Tip{read}{appc}{.2}{7.7}{12}{JSON読取の所在}{\C{environment.storage.read_json}を直接使い、weatherの私有名へ横断しない。旧\C{cli.read_json}のimportは同じ実体を再exportする。標準JSON以外の暗黙補正をしない。}
\Tip{exit}{outc}{14.1}{7.7}{11.7}{利用者へ返す意味}{0：操作成功、またはlanded。\par 2：stoppedの部分結果を保存。\par 1：入力・実行・保存などのエラー。\par argparseのusage終了2は別。操作例と失敗後の判断はCOMMANDSへ。}
\end{Map}

\Page{toolc}{PROGRAM-TESTS}{検査と道具：何を守るために存在するか}{横断枝 / tests と tools → 対象契約と本体への接続}{試験数や図の存在だけで構造を受け入れない。守る意味と、その証拠が届く範囲を分ける。}
\begin{Map}
\Card{weather}{weatherc}{.2}{.3}{7.2}{\C{test_flight_weather.py}}{取得要求、raw、支持/補間/地表。\par \C{test_gfs_acquisition.py}：共有間隔・通信制限・resume。\par \C{test_gfs_raw_reuse.py}：候補照合・独立copy・残量/破損を注入検査。}
\Card{flight}{dync}{9.4}{.3}{7.2}{\C{test_trajectory.py}}{解析場・独立式・破裂/着地/停止。\par flightの変更を守る。\par \C{test_flight_numerics.py}はSciPy接続と根を独立検査。}
\Card{cli}{outc}{18.6}{.3}{7.2}{集合の純粋な意味}{\C{tests/test_ensemble.py}\par 同じ抽選・方式への展開、点/線/面、経験順位、相/終端・選択母数。\par 核を試し、実UIや科学精度の受入とは分ける。}
\hypertarget{PROGRAM-TOOLS}{}
\File{fixture}{weatherc}{.2}{4.7}{7.2}{保存標本とJRA接続}{\C{make_flight_fixture.py}：保存場から小標本を切り出す。\par \C{prepare_jra3q_flight_fixture.py}：固定原本を変換。\par \C{test_jra3q_flight_field.py}：列場・製品契約・保存振分け・飛行停止。\par \C{test_jra3q_surface_field.py}：Bの順序・固定join・支持/欠測・別schema。}
\File{proto}{toolc}{9.4}{4.7}{7.2}{backend・frontendの試験}{\C{backend/tests/test_ensemble.py}：固定計画・全台帳・取消/再開・保存。\par \C{frontend/src/ensemble.test.ts}：参照・描画adapterの契約。\par 風統計は\C{test_climate*.py / climate.test.ts}。\par \C{test_weather_raw_reuse.py}：登録完成assetからの取得。\C{restoration.test.ts}：未読参照・寿命/保存境界。}
\File{manage}{toolc}{18.6}{4.7}{7.2}{管理・生成}{check_state/health/contract\par run_checks / prepare_workspace\par build_docs / build_guides\par 計算の実行依存ではない。}
\Tip{judge}{ink}{.2}{9.3}{25.6}{変更後は親の用途へ戻る}{旧GFS/JRA等の試作は比較根拠として保持する。\C{test_flight_cli.py}の実場再現・出力署名、backendの実worker試験、画面の時刻/相/選択試験は継続する。これらと、原一覧/GRIBを実取得して再閲覧する証拠、資料の通読、科学的な精度評価は別である。取得の一成功から全国・全期間・全モデルの受入へ広げない。}
\end{Map}

\Page{ink}{PROGRAM-CHANGE}{変更したいことから、読む枝を決める}{逆引き / 最初の担当 → 横の契約 → 利用と試験への影響}{この図で終わらず、該当枝へ降り、最後に全体の目的へ戻る。配置が見つかることと、変更案が妥当なことは別である。}
\begin{Map}
\Card{product}{weatherc}{.2}{.3}{12}{気象の取得・意味を変える}{操作/保存：\Go{WEB-PREPARATION}{取得UI} → \Go{WEB-WEATHER}{backend/weather}\par GFS通信/製品：nomads.py、\Go{WEATHER-ACQUIRE}{gfs.py / gfs_contract.py}\par 保存schema分岐：storage。GFS再構成：\Go{WEATHER-FIELD}{bundle / fields}\par JRA製品/strict：\Go{WEATHER-MODEL}{jra3q / model_levels}\par 暫定B再構成：\Go{WEATHER-SURFACE}{model_surface / 明示policy}\par 支持/品質と全飛行へ影響。根拠はWEATHER / ENV / THEORY。}
\Card{model}{dync}{14.1}{.3}{11.7}{鉛直則・状態量・破裂を変える}{最初：\Go{FLIGHT-MODELS}{models.py / events.py} と config.py\par 次：\Go{DYNAMICS-SIMULATE}{trajectory.py} の状態・要求量・記録\par 影響：結果schema・停止・比較試験\par 根拠：THEORY / ENV / 実装ノート}
\Card{solver}{dync}{.2}{5}{12}{積分器を変える}{最初：\Go{FLIGHT-NUMERICS}{numerics.py} の呼出契約\par 次：trajectoryの相・支持・event・出力点\par 影響：依存/許容誤差/来歴/解析解との比較\par 採用点/誤差の意味を保ち、一般法はSciPyへ委譲。}
\Card{output}{outc}{14.1}{5}{11.7}{操作・結果表示を変える}{集合：\Go{WEB-ENSEMBLE-FRONTEND}{画面} → \Go{WEB-ENSEMBLE-BACKEND}{台帳/保存} → \Go{ENSEMBLE-CORE}{抽選/集計}\par CLI単独：\Go{PROGRAM-CLI}{cli.py} / \Go{CLI-EXPORT}{results/}\par 風統計：\Go{WEB-CLIMATE-FRONTEND}{資料adapter/共通図} → \Go{WEB-CLIMATE-BACKEND}{climate/}\par 次：草案/固定結果・API・保存版・再閲覧\par 計算の式を表示層へ複写しない。\par 根拠：COMMANDS / 実装ノート / CLI試験}
\Tip{end}{ink}{.2}{10}{25.6}{この地図から分かることと、残ること}{実パッケージの責務、主要依存、変更の最初の場所と影響先を辿れる。全組合せの物理的採用、全国全期間の取得、実飛行の予測精度を認定する資料ではない。読めない枝が増えたら、図を縮めて押し込む前にコード側の責務を再検討する。}
\end{Map}
% LOGIC-END:PROGRAM-CORE
\end{document}
